This chapter develops several canonical examples of sample spaces and elementary probability from the ground up. The emphasis is on understanding the experiment, deciding what should count as an elementary outcome, constructing the sample space, defining events, and assigning probabilities only after the mathematical model has been made explicit.

Different representations of the same probabilistic structure should be used throughout the chapter: lists, tables, trees, diagrams, and symbolic notation. The goal is to show that a sample space is an abstract object produced by analysis and modelling, not a collection of observed data.
